\documentclass[../main]{subfiles}
\begin{document}

\section{Teorema de Bernoulli}

El teorema de Bernoulli es un principio fundamental en la mecánica de fluidos que describe el comportamiento de un fluido en movimiento. Este teorema es ampliamente utilizado en diversas aplicaciones, desde la aerodinámica hasta la hidráulica, y proporciona una relación entre la velocidad, la presión y la altura de un fluido en un flujo continuo.

\subsection{Fundamentos Teóricos}

El teorema de Bernoulli establece que, en un flujo de fluido sin fricción y sin viscosidad, la suma de la presión hidrostática, la energía cinética y la energía potencial por unidad de volumen permanece constante a lo largo de una línea de corriente. Matemáticamente, esto se expresa como:

\info{Teorema de Bernoulli}{
El teorema de Bernoulli se puede expresar de la siguiente manera:
\begin{equation}
P + \frac{1}{2}\rho v^2 + \rho gh = \text{constante}
\end{equation}
donde:
\begin{itemize}
    \item $P$ es la presión del fluido en el punto considerado (en Pa).
    \item $\rho$ es la densidad del fluido (en kg/m$^3$).
    \item $v$ es la velocidad del fluido en el punto considerado (en m/s).
    \item $g$ es la aceleración debida a la gravedad (en m/s$^2$).
    \item $h$ es la altura del fluido sobre un punto de referencia (en m).
\end{itemize}
}

El teorema de Bernoulli se deriva a partir del principio de conservación de la energía, aplicado a un sistema fluido en movimiento en un campo gravitatorio constante.

\subsection{Aplicaciones}

El teorema de Bernoulli tiene una amplia gama de aplicaciones en la ingeniería y la física, incluyendo:

\begin{itemize}
    \item \textbf{Aerodinámica}: Se utiliza para analizar el flujo de aire alrededor de objetos en movimiento, como aviones, para determinar la generación de sustentación y resistencia.
    \item \textbf{Hidráulica}: Se aplica para estudiar el flujo de agua en tuberías, canales y sistemas de riego, para el diseño de sistemas de distribución de agua y estructuras hidráulicas.
    \item \textbf{Medicina}: Se emplea para entender el flujo sanguíneo en el sistema circulatorio humano y para analizar el funcionamiento de dispositivos médicos como catéteres y sondas.
    \item \textbf{Meteorología}: Se utiliza para modelar el flujo de aire en la atmósfera terrestre y para predecir el clima y los patrones meteorológicos.
    \item \textbf{Ingeniería de Ventilación}: Se aplica para diseñar sistemas de ventilación en edificios y vehículos para garantizar una distribución adecuada del aire.
\end{itemize}

\subsection{Preguntas para Investigar}

\begin{enumerate}
    \item ¿Qué condiciones deben cumplirse para que el teorema de Bernoulli sea aplicable?
    \item ¿Cómo se puede utilizar el teorema de Bernoulli para explicar el fenómeno del efecto Venturi?
    \item ¿Cuál es la diferencia entre flujo laminar y flujo turbulento en términos del teorema de Bernoulli?
    \item ¿Cómo se relaciona el teorema de Bernoulli con el principio de conservación de la energía?
    \item ¿Qué efecto tiene la viscosidad del fluido en la aplicación del teorema de Bernoulli?
    \item ¿Cómo se puede aplicar el teorema de Bernoulli en el diseño de sistemas de riego por goteo?
    \item ¿Cuál es la importancia del teorema de Bernoulli en la industria aeronáutica?
    \item ¿Qué sucede con la presión de un fluido cuando su velocidad aumenta de acuerdo con el teorema de Bernoulli?
    \item ¿Cómo se puede utilizar el teorema de Bernoulli para calcular la velocidad del flujo en una tubería?
    \item ¿Qué implicaciones tiene el teorema de Bernoulli en el diseño de sistemas de calefacción y aire acondicionado?
\end{enumerate}
\subsection{Problemas Prácticos}

\begin{enumerate}
    \item Un tanque de agua tiene una altura de 10 metros. Si se abre una válvula en la parte inferior del tanque, ¿cuál es la velocidad del agua en la salida?
    \item Un avión vuela a una altitud de 8000 metros y a una velocidad de 200 m/s. ¿Cuál es la presión del aire en el exterior del avión, si la presión atmosférica al nivel del mar es de 101325 Pa?
    \item Una manguera de jardín tiene un diámetro de 2 cm y suministra agua a una velocidad de 4 m/s. ¿Cuál es la presión del agua en el extremo de la manguera, si la manguera se encuentra a una altura de 1 metro por encima del suelo?
    \item Un conducto tiene una sección transversal de 0.5 m\textsuperscript{2} y una velocidad de agua de 6 m/s. Si la presión del agua en el conducto es de 50000 Pa, ¿cuál es la altura del agua sobre un punto de referencia?
    \item Un submarino se sumerge a una profundidad de 200 metros en el océano. Si la presión atmosférica al nivel del mar es de 101325 Pa, ¿cuál es la presión del agua en el exterior del submarino?
    \item Una fuente de agua a una altura de 15 metros suministra agua a través de una tubería de 5 cm de diámetro. Si la velocidad del agua en la tubería es de 8 m/s, ¿cuál es la presión del agua en la tubería en Pa?
    \item Se suelta un globo lleno de aire desde el suelo. Si el globo alcanza una altitud de 500 metros, ¿cuál es la presión del aire en el interior del globo en Pa?
    \item Un ventilador expulsa aire a una velocidad de 10 m/s a través de una sección transversal de 0.2 m\textsuperscript{2}. Si la presión del aire en el ventilador es de 10000 Pa, ¿cuál es la potencia del ventilador en W?
    \item Una cascada tiene una altura de 20 metros. Si el flujo de agua tiene una velocidad de 5 m/s en la parte superior de la cascada, ¿cuál es la velocidad del agua en la base de la cascada?
    \item Se instala un sistema de riego por aspersión en un jardín. Si la presión del agua en la tubería es de 30000 Pa y la altura de la tubería es de 2 metros por encima del suelo, ¿cuál es la velocidad del agua en las boquillas de pulverización?
\end{enumerate}

\end{document}
